\subsection{Time delay embedding}
It is the technique for reconstructing the underlying dynamics of the time series by transforming a signal into a higher order dimensional state space representation. The time delay embedding vector is defined as: 
\begin{equation}
\mathbf{z}_t = \begin{bmatrix}
x_t \\
x_{t-\tau} \\
x_{t-2\tau} \\
\vdots \\
x_{t-(d-1)\tau}
\end{bmatrix} 
\label{eq:tde_definition}
\end{equation}
In a sliding window implantation of length $W$(the buffer window size),the time-delay embedding vector $x_t \in R^d$ of dimension $d$ is extracted from the last element of the current window. Because the index offsets defining each embedding dimensions remain constant across iterations, they can be pre-calculated during initialization according to:
\begin{equation}
\text{index}_i = (W - 1) - i \cdot \tau, \quad i = 0, 1, \ldots, d-1
\label{eq:tde_indices}
\end{equation}
The values from the sliding window into the time delay embedding vector get mapped with this $index_i$ mapping. 